% Preámbulo común para las soluciones de los prácticos de Programación 3
\documentclass[11pt,a4paper]{article}

\usepackage[utf8]{inputenc}
\usepackage[T1]{fontenc}
\usepackage{lmodern}
\usepackage[spanish,es-noshorthands]{babel}
\usepackage[margin=2.3cm]{geometry}
\usepackage{amsmath,amssymb,amsthm}
\usepackage{enumitem}
\usepackage{xcolor}
\usepackage{tikz}
\usetikzlibrary{arrows.meta,positioning,shapes,calc}
\usepackage[noend]{algpseudocode}
\usepackage{algorithm}
\usepackage{booktabs}
\usepackage{tcolorbox}
\usepackage[hidelinks]{hyperref}

\floatname{algorithm}{Algoritmo}
\renewcommand{\algorithmicrequire}{\textbf{Entrada:}}
\renewcommand{\algorithmicensure}{\textbf{Salida:}}
\algrenewcommand\algorithmicwhile{\textbf{mientras}}
\algrenewcommand\algorithmicdo{\textbf{hacer}}
\algrenewcommand\algorithmicfor{\textbf{para}}
\algrenewcommand\algorithmicforall{\textbf{para todo}}
\algrenewcommand\algorithmicif{\textbf{si}}
\algrenewcommand\algorithmicthen{\textbf{entonces}}
\algrenewcommand\algorithmicelse{\textbf{sino}}
\algrenewcommand\algorithmicreturn{\textbf{retornar}}
\algrenewcommand\algorithmicfunction{\textbf{función}}
\algrenewcommand\algorithmicprocedure{\textbf{procedimiento}}

\theoremstyle{plain}
\newtheorem{lema}{Lema}[section]
\newtheorem{teorema}[lema]{Teorema}
\newtheorem{proposicion}[lema]{Proposición}
\theoremstyle{definition}
\newtheorem{definicion}[lema]{Definición}
\theoremstyle{remark}
\newtheorem*{observacion}{Observación}

% Caja para el enunciado de cada ejercicio
\newtcolorbox{enunciado}{colback=gray!6,colframe=gray!50,boxrule=0.4pt,arc=2pt,left=6pt,right=6pt,top=4pt,bottom=4pt,fontupper=\small}
% Nota que vincula el ejercicio con parciales/exámenes
\newcommand{\examref}[1]{\par\noindent{\small\color{blue!50!black}\textbf{Relación con parciales/exámenes:} #1}\par\medskip}

\newcommand{\NP}{\mathsf{NP}}
\newcommand{\Pclass}{\mathsf{P}}
\newcommand{\Oh}{\mathcal{O}}

\setlength{\parskip}{0.4em}
\setlength{\parindent}{0pt}


\newcommand{\Equiv}{\textsc{Equiv}}
\newcommand{\nil}{\bot}
\newcommand{\ind}[1]{\mathbf{1}_{\{#1\}}}

\title{Práctico 7 --- Dividir y Conquistar}
\author{Programación 3 --- FING, UdelaR}
\date{}

\begin{document}
\maketitle

\section*{Resultados del libro que se usan}

En todo el práctico usamos los resultados de la sección 5.1--5.2 de Kleinberg y Tardos,
\emph{Algorithm Design}. Los enunciamos acá una sola vez.

\begin{teorema}[K\&T (5.1) y (5.2)]\label{t:52}
Sea $T$ una función que cumple, para alguna constante $c$,
\[
T(n)\le 2\,T(n/2)+cn \quad (n>2),\qquad T(2)\le c. \tag{5.1}
\]
Entonces $T(n)$ es $\Oh(n\log n)$ para $n>1$ \emph{(5.2)}.
\end{teorema}

\begin{teorema}[K\&T (5.3), (5.4), (5.5)]\label{t:q}
Sea $T$ tal que $T(n)\le q\,T(n/2)+cn$ para $n>2$ y $T(2)\le c$ \emph{(5.3)}. Si $q>2$
entonces $T(n)$ es $\Oh(n^{\log_2 q})$ \emph{(5.4)}; si $q=1$ entonces $T(n)$ es
$\Oh(n)$ \emph{(5.5)}.
\end{teorema}

En el libro estos resultados se prueban para $n$ potencia de $2$, desenrollando la
recurrencia o con el árbol de recursión: en el nivel $j$ del árbol hay $2^j$
subproblemas de tamaño $n/2^j$, cada uno aporta a lo sumo $c\,n/2^j$ de trabajo no
recursivo, así que cada nivel aporta a lo sumo $cn$; hay $\log_2 n$ niveles y el total es
a lo sumo $cn\log_2 n$. Cuando $n$ no es potencia de $2$, en cada ejercicio indicamos
cómo se resuelve ese detalle.

%======================================================================
\section{Ejercicio 1}

\begin{enunciado}
\textbf{(Kleinberg \& Tardos, Ex.\ 5.3).} Queremos analizar $n$ tarjetas bancarias que
fueron confiscadas por sospecha de fraude. Cada tarjeta está asociada a una única cuenta
de banco, aunque cada cuenta puede estar asociada a más de una tarjeta; dos tarjetas son
equivalentes si están asociadas a una misma cuenta.

Leer y comparar los datos de cada tarjeta es un proceso costoso, pero el banco cuenta con
un dispositivo al que se ingresan dos tarjetas y responde rápidamente si son equivalentes
o no.

Nos interesa saber si entre las $n$ tarjetas confiscadas hay algún conjunto de más de
$n/2$ tarjetas que son todas equivalentes entre sí.
\begin{enumerate}[label=(\alph*)]
\item Dé un algoritmo de tipo divide y vencerás para solucionar el problema realizando
$\Oh(n\log n)$ pruebas de equivalencia con el dispositivo mencionado.
\item Demuestre la corrección de su algoritmo.
\item Demuestre que la cantidad de pruebas de equivalencia que realiza su algoritmo es
$\Oh(n\log n)$ como se pide. Enuncie los resultados teóricos del libro de referencia que
utilice.
\end{enumerate}
\emph{Sugerencia:} Note que no puede haber más de una clase de equivalencia (conjunto de
tarjetas equivalentes entre sí) de tamaño mayor a $n/2$. Utilice este hecho para definir
un algoritmo que devuelve, si existe, un representante de dicha clase de equivalencia.
\end{enunciado}

\paragraph{Notación.} Para una secuencia (multiconjunto) de tarjetas $S$ con $|S|=n$,
la relación ``estar asociadas a la misma cuenta'' es de equivalencia, y
particiona $S$ en clases. Decimos que una clase $C$ es \emph{mayoritaria en $S$} si
$|C\cap S|>|S|/2$. El dispositivo es la función $\Equiv(x,y)$, que devuelve
\textbf{verdadero} si y solo si $x$ e $y$ son equivalentes; lo único que contamos es la
cantidad de invocaciones a $\Equiv$.

\begin{lema}[Unicidad]\label{l:unica}
En una secuencia $S$ hay a lo sumo una clase mayoritaria.
\end{lema}
\begin{proof}
Si $C\ne C'$ fueran dos clases mayoritarias, al ser clases distintas son disjuntas, y
entonces $|S|\ge|C\cap S|+|C'\cap S|>|S|/2+|S|/2=|S|$, absurdo.
\end{proof}

\begin{lema}[Herencia]\label{l:herencia}
Sea $S$ partida en dos secuencias $S_1$, $S_2$ (o sea, $|S|=|S_1|+|S_2|$ y cada tarjeta
de $S$ está exactamente en una de las dos). Si $C$ es mayoritaria en $S$, entonces $C$ es
mayoritaria en $S_1$ o en $S_2$.
\end{lema}
\begin{proof}
Si no lo fuera en ninguna, $|C\cap S_1|\le|S_1|/2$ y $|C\cap S_2|\le|S_2|/2$, y sumando
$|C\cap S|\le(|S_1|+|S_2|)/2=|S|/2$, lo que contradice que $C$ es mayoritaria en $S$.
\end{proof}

Ojo: la recíproca no vale. Una clase puede ser mayoritaria en una mitad y no en $S$; por
eso el algoritmo tiene que \emph{verificar} a cada candidato contra toda la secuencia.

\subsection*{(a) Algoritmo}

La función $\textsc{Representante}(S)$ devuelve una tarjeta de la clase mayoritaria de $S$
si esa clase existe, y $\nil$ si no existe. Para resolver el problema se invoca
$\textsc{Representante}(s_1,\dots,s_n)$: hay más de $n/2$ tarjetas equivalentes entre sí
si y solo si el resultado es distinto de $\nil$.

\begin{algorithm}[H]
\caption{Representante de la clase mayoritaria}
\begin{algorithmic}[1]
\Require Secuencia de tarjetas $S=(s_1,\dots,s_n)$, $n\ge 1$
\Ensure Una tarjeta de la clase mayoritaria de $S$, o $\nil$ si no hay clase mayoritaria
\Function{Representante}{$S$}
  \If{$n=1$}
    \State \Return $s_1$ \Comment{caso base: la única clase tiene $1>1/2$ elementos}
  \EndIf
  \State $m\gets\lfloor n/2\rfloor$
  \State $S_1\gets(s_1,\dots,s_m)$;\quad $S_2\gets(s_{m+1},\dots,s_n)$ \Comment{dividir}
  \State $r_1\gets\Call{Representante}{S_1}$;\quad $r_2\gets\Call{Representante}{S_2}$
  \For{cada $r\in\{r_1,r_2\}$ con $r\neq\nil$} \Comment{combinar: verificar candidatos}
    \State $cnt\gets 0$
    \For{$k\gets 1$ \textbf{a} $n$}
      \If{$\Call{Equiv}{r,s_k}$}
        \State $cnt\gets cnt+1$
      \EndIf
    \EndFor
    \If{$cnt>n/2$}
      \State \Return $r$
    \EndIf
  \EndFor
  \State \Return $\nil$
\EndFunction
\end{algorithmic}
\end{algorithm}

En la línea del \textbf{para} interno se puede evitar comparar $r$ consigo mismo (se
cuenta $1$ sin usar el dispositivo); no cambia el orden.

\subsection*{(b) Correctitud}

\begin{proposicion}
Para toda secuencia $S$ con $n\ge1$ tarjetas, $\textsc{Representante}(S)$ termina, y
devuelve una tarjeta de la clase mayoritaria de $S$ si existe; si no existe, devuelve
$\nil$.
\end{proposicion}
\begin{proof}
Por inducción completa en $n$.

\emph{Terminación.} Si $n\ge 2$ entonces $1\le m=\lfloor n/2\rfloor<n$ y
$1\le n-m<n$, así que las llamadas recursivas son sobre secuencias no vacías y
estrictamente más chicas; los ciclos son acotados.

\emph{Caso base, $n=1$.} La única clase de $S$ es la de $s_1$, con $1>1/2$ elementos, así
que es mayoritaria; el algoritmo devuelve $s_1$, que es de esa clase.

\emph{Paso inductivo, $n\ge2$.} Suponemos (hipótesis inductiva, HI) que la proposición
vale para toda secuencia de tamaño entre $1$ y $n-1$, en particular para $S_1$ y $S_2$.
Observemos primero que el ciclo interno calcula, en $cnt$, exactamente
$|[r]\cap S|$, la cantidad de tarjetas de $S$ en la clase $[r]$ de $r$. Por lo tanto el
algoritmo devuelve un candidato $r$ si y solo si $[r]$ es mayoritaria en $S$. Hay dos
casos.

\begin{itemize}
\item \emph{$S$ tiene una clase mayoritaria $C$.} Por el Lema~\ref{l:herencia}, $C$ es
mayoritaria en $S_1$ o en $S_2$; digamos en $S_i$. Por el Lema~\ref{l:unica}, $C$ es
\emph{la} clase mayoritaria de $S_i$, y por la HI $r_i\neq\nil$ y $r_i\in C$. Entonces, al
verificar $r_i$, se obtiene $cnt=|C\cap S|>n/2$ y se devuelve $r_i\in C$, salvo que antes
se haya devuelto el otro candidato $r_j$. Pero eso solo pasa si $[r_j]$ es mayoritaria en
$S$, y por el Lema~\ref{l:unica} $[r_j]=C$. En cualquier caso se devuelve una tarjeta de
$C$.
\item \emph{$S$ no tiene clase mayoritaria.} Ningún candidato $r$ puede cumplir
$cnt=|[r]\cap S|>n/2$, porque si no $[r]$ sería mayoritaria. Entonces no se ejecuta
ningún \textbf{retornar} $r$ y el algoritmo devuelve $\nil$. \qedhere
\end{itemize}
\end{proof}

En particular, como la clase mayoritaria de las $n$ tarjetas es justamente un conjunto
de más de $n/2$ tarjetas equivalentes entre sí (y cualquier conjunto así está contenido
en una única clase, que entonces es mayoritaria), la respuesta al problema es
``sí'' si y solo si $\textsc{Representante}(s_1,\dots,s_n)\ne\nil$.

\subsection*{(c) Cantidad de pruebas de equivalencia}

Sea $T(n)$ la máxima cantidad de invocaciones a $\Equiv$ que hace
$\textsc{Representante}$ sobre una secuencia de $n$ tarjetas.
\begin{itemize}
\item Caso base: con $n=1$ no se usa el dispositivo, así que $T(1)=0$.
\item Caso recursivo: dividir no usa el dispositivo; las dos llamadas recursivas hacen a
lo sumo $T(\lfloor n/2\rfloor)+T(\lceil n/2\rceil)$ pruebas; la verificación hace a lo
sumo $n$ pruebas por candidato y hay a lo sumo $2$ candidatos, o sea a lo sumo $2n$.
\end{itemize}
Entonces
\[
T(1)=0,\qquad T(n)\le T(\lfloor n/2\rfloor)+T(\lceil n/2\rceil)+2n\quad(n\ge2).
\]

\paragraph{Con el resultado del libro.} Si $n$ es potencia de $2$ las dos mitades tienen
tamaño $n/2$ y queda $T(n)\le 2T(n/2)+2n\le 2T(n/2)+4n$ para $n>2$, con
$T(2)\le 2T(1)+4=4$. Esta es la recurrencia (5.1) con $c=4$, y por (5.2)
(Teorema~\ref{t:52}) $T(n)$ es $\Oh(n\log n)$.

\paragraph{Para $n$ arbitrario.} Probamos directamente, por inducción completa, que
$T(n)\le 2n\lceil\log_2 n\rceil$ para todo $n\ge1$. Para $n=1$ ambos lados valen $0$. Para
$n\ge2$ usamos que $\lceil\log_2\lceil n/2\rceil\rceil=\lceil\log_2 n\rceil-1$ (si
$2^{k-1}<n\le 2^k$ entonces $2^{k-2}<\lceil n/2\rceil\le 2^{k-1}$) y que, como
$1\le\lfloor n/2\rfloor\le\lceil n/2\rceil$ y $\log_2$ es creciente, también
$\lceil\log_2\lfloor n/2\rfloor\rceil\le\lceil\log_2 n\rceil-1$:
\[
T(n)\le 2\lfloor n/2\rfloor(\lceil\log_2 n\rceil-1)+2\lceil n/2\rceil(\lceil\log_2
n\rceil-1)+2n = 2n(\lceil\log_2 n\rceil-1)+2n=2n\lceil\log_2 n\rceil .
\]
Así que la cantidad de pruebas de equivalencia es $\Oh(n\log n)$, como se pedía.

El resto del trabajo (armar $S_1$ y $S_2$ y llevar los contadores) es proporcional a la
cantidad de pruebas en cada nivel, así que el tiempo total también es $\Oh(n\log n)$.

\begin{observacion}
Existe incluso una solución con $\Oh(n)$ pruebas (aparear las
tarjetas, descartar los pares no equivalentes y quedarse con una de cada par equivalente,
lo que da $T(n)\le T(n/2)+\Oh(n)$, que es $\Oh(n)$ por (5.5)). No se pide acá.
\end{observacion}

%======================================================================
\section{Ejercicio 2}

\begin{enunciado}
\textbf{(Kleinberg \& Tardos, Ex.\ 5.2).} Sea $L_1,L_2,\dots,L_n$ una secuencia de
números. Decimos que un par de índices $i,j$, $1\le i<j\le n$, es una \emph{inversión
significativa} si se cumple $L_i>2L_j$.
\begin{enumerate}[label=(\alph*)]
\item Dé un algoritmo para calcular la cantidad de inversiones significativas. Su
algoritmo debe admitir una implementación con tiempo de ejecución $\Oh(n\log n)$.
\emph{Sugerencia:} Adapte la función Merge-and-count de la sección 5.3 del libro; observe
que el orden relativo de los elementos de una secuencia no se altera si todos se
multiplican por dos.
\item Demuestre la corrección de su algoritmo. \emph{Sugerencia:} Demuestre que para $n$
mayor que una constante (determinada por la condición del paso base de su algoritmo), la
cantidad de inversiones en $L$, $R(L)$, se puede expresar como
\[ R(L)=R(A)+r+R(B), \tag{2.1}\]
donde $A$ y $B$ son las partes en que se subdivide el problema original, de tamaños $n_A$
y $n_B$, respectivamente, y
\[ r=\sum_{j=1}^{n_B}\sum_{i=1}^{n_A}\ind{A_i>2B_j}. \tag{2.2}\]
Luego demuestre por inducción que su algoritmo calcula $R(L)$.
\item Demuestre que su algoritmo admite una implementación con tiempo de ejecución que
es $\Oh(n\log n)$. Enuncie los resultados teóricos del libro de referencia que utilice.
\end{enumerate}
\end{enunciado}

\examref{Examen julio 2022, Ej.~1 (contar inversiones ``comunes'' con una variante de
MergeSort, Merge-and-Count / Sort-and-Count de la sección 5.3, y justificar
$\Oh(n\log n)$ con la recurrencia (5.1)). Este ejercicio es la misma idea, cambiando la
condición $L_i>L_j$ por $L_i>2L_j$; la diferencia es que el conteo ya no se puede hacer
dentro del mismo intercalado, y conviene hacerlo en una pasada aparte.}

\paragraph{Notación.} Para una secuencia $L=(L_1,\dots,L_n)$ sea
$I(L)=\{(i,j): 1\le i<j\le n,\ L_i>2L_j\}$ y $R(L)=|I(L)|$. Para dos secuencias $A$ (de
largo $n_A$) y $B$ (de largo $n_B$) definimos
\[
r(A,B)=\sum_{j=1}^{n_B}\sum_{i=1}^{n_A}\ind{A_i>2B_j}
=\bigl|\{(i,j): 1\le i\le n_A,\ 1\le j\le n_B,\ A_i>2B_j\}\bigr|,
\]
la cantidad de ``inversiones significativas cruzadas'' entre $A$ y $B$.

\subsection*{(a) Algoritmo}

Igual que en Sort-and-Count, la función recursiva devuelve un par: la cantidad de
inversiones significativas y la secuencia ordenada. Ordenar sirve para que el conteo de
inversiones cruzadas se pueda hacer en tiempo lineal.

\begin{algorithm}[H]
\caption{Conteo de inversiones significativas cruzadas}
\begin{algorithmic}[1]
\Require $A=(a_1,\dots,a_p)$ y $B=(b_1,\dots,b_q)$, ambas ordenadas en forma creciente
\Ensure $r(A,B)=|\{(i,j): a_i>2b_j\}|$
\Function{ContarCruzadas}{$A,B$}
  \State $cnt\gets 0$;\quad $i\gets 1$
  \For{$j\gets 1$ \textbf{a} $q$}
    \While{$i\le p$ \textbf{y} $a_i\le 2b_j$}
      \State $i\gets i+1$
    \EndWhile
    \State $cnt\gets cnt+(p-i+1)$ \Comment{$a_i,\dots,a_p$ son exactamente los $>2b_j$}
  \EndFor
  \State \Return $cnt$
\EndFunction
\end{algorithmic}
\end{algorithm}

\begin{algorithm}[H]
\caption{Ordenar y contar inversiones significativas}
\begin{algorithmic}[1]
\Require Secuencia $L=(L_1,\dots,L_n)$, $n\ge1$
\Ensure El par $(R(L),L')$, con $L'$ la secuencia $L$ ordenada en forma creciente
\Function{OrdenarYContar}{$L$}
  \If{$n=1$}
    \State \Return $(0,L)$
  \EndIf
  \State $m\gets\lfloor n/2\rfloor$;\quad $A\gets(L_1,\dots,L_m)$;\quad
         $B\gets(L_{m+1},\dots,L_n)$ \Comment{dividir}
  \State $(r_A,A')\gets\Call{OrdenarYContar}{A}$
  \State $(r_B,B')\gets\Call{OrdenarYContar}{B}$
  \State $r\gets\Call{ContarCruzadas}{A',B'}$ \Comment{combinar: contar}
  \State $L'\gets\Call{Merge}{A',B'}$ \Comment{combinar: intercalado usual de MergeSort}
  \State \Return $(r_A+r+r_B,\ L')$
\EndFunction
\end{algorithmic}
\end{algorithm}

\textsc{Merge} es el intercalado de MergeSort: recorre $A'$ y $B'$ con dos cursores y en
cada paso agrega al final de la salida el menor de los dos elementos apuntados; termina
copiando el resto de la secuencia que no se agotó.

\paragraph{Por qué no alcanza con modificar el intercalado directamente.} En
Merge-and-Count, cuando se agrega $b_j$ a la salida se suman los elementos de $A$ que
quedan, porque son justo los $>b_j$. Acá necesitamos los $>2b_j$, que son un sufijo de $A'$
distinto (en general más corto). Por eso usamos un segundo cursor en una pasada propia,
\textsc{ContarCruzadas}. Esto es lo que dice la sugerencia: como $x\mapsto 2x$ es
creciente, la secuencia $2b_1\le 2b_2\le\dots\le 2b_q$ también está ordenada, y se puede
``intercalar'' $A'$ con $2B'$ para contar.

\paragraph{Ejemplo.} $L=(5,1,7,2)$: las inversiones significativas son $(1,2)$
($5>2$), $(1,4)$ ($5>4$) y $(3,4)$ ($7>4$); $R(L)=3$. El algoritmo calcula
$R((5,1))=1$, $R((7,2))=1$, y con $A'=(1,5)$, $B'=(2,7)$: para $b_1=2$ quedan los $a_i>4$,
o sea $\{5\}$, que aporta $1$; para $b_2=7$ no queda ningún $a_i>14$. Total
$1+1+1=3$.

\subsection*{(b) Correctitud}

\begin{lema}[Descomposición (2.1)--(2.2)]\label{l:desc}
Sea $n\ge 2$, $m=\lfloor n/2\rfloor$, $A=(L_1,\dots,L_m)$ y $B=(L_{m+1},\dots,L_n)$ (o
sea $A_i=L_i$ y $B_j=L_{m+j}$, con $n_A=m$, $n_B=n-m$). Entonces
$R(L)=R(A)+r(A,B)+R(B)$.
\end{lema}
\begin{proof}
Partimos $I(L)$ en tres conjuntos disjuntos, según dónde caen $i$ y $j$ (recordar que
$i<j$):
\begin{itemize}
\item $I_1=\{(i,j)\in I(L): j\le m\}$. Ambos índices están en $A$ y la condición es
$A_i>2A_j$ con $i<j\le m$; o sea $I_1=I(A)$ y $|I_1|=R(A)$.
\item $I_3=\{(i,j)\in I(L): i> m\}$. Con el cambio $i=m+i'$, $j=m+j'$, la condición queda
$B_{i'}>2B_{j'}$ con $1\le i'<j'\le n_B$; es una biyección con $I(B)$, así que
$|I_3|=R(B)$.
\item $I_2=\{(i,j)\in I(L): i\le m<j\}$. Acá la condición $i<j$ se cumple siempre, y con
$j=m+j'$ la condición es $A_i>2B_{j'}$ con $1\le i\le n_A$, $1\le j'\le n_B$. Así que
$|I_2|=\sum_{j'=1}^{n_B}\sum_{i=1}^{n_A}\ind{A_i>2B_{j'}}=r(A,B)$.
\end{itemize}
Como todo par de $I(L)$ cae en exactamente uno de los tres casos,
$R(L)=|I_1|+|I_2|+|I_3|=R(A)+r(A,B)+R(B)$.
\end{proof}

\begin{lema}[Invariancia de $r$]\label{l:inv}
Si $A'$ es una permutación de $A$ y $B'$ es una permutación de $B$, entonces
$r(A',B')=r(A,B)$.
\end{lema}
\begin{proof}
$r(A,B)$ cuenta todos los pares (elemento de $A$, elemento de $B$) con el primero mayor
que el doble del segundo, sin mirar posiciones. Si $A'_i=A_{\sigma(i)}$ y
$B'_j=B_{\tau(j)}$ con $\sigma,\tau$ biyecciones, la suma doble para $(A',B')$ tiene
exactamente los mismos términos que la de $(A,B)$, solo que reordenados.
\end{proof}

\begin{lema}[Correctitud de \textsc{ContarCruzadas}]\label{l:cc}
Si $A=(a_1,\dots,a_p)$ y $B=(b_1,\dots,b_q)$ están ordenadas en forma creciente,
$\textsc{ContarCruzadas}(A,B)$ devuelve $r(A,B)$.
\end{lema}
\begin{proof}
Sea $k_j=|\{i: a_i\le 2b_j\}|$. Como $A$ es creciente, los $a_i\le 2b_j$ son exactamente
$a_1,\dots,a_{k_j}$ (un prefijo) y los $a_i>2b_j$ son $a_{k_j+1},\dots,a_p$, que son
$p-k_j$. Como $B$ es creciente, $2b_{j-1}\le 2b_j$ y por lo tanto
$k_{j-1}\le k_j$.

Probamos por inducción en $j$ el invariante: \emph{al terminar la iteración $j$ del
\textbf{para} se cumple $i=k_j+1$ y $cnt=\sum_{j'=1}^{j}(p-k_{j'})$}. Al comenzar la
iteración $j$ vale $i=k_{j-1}+1$ (con $k_0=0$ para $j=1$). Todos los $a_{i'}$ con
$i'\le k_j$ cumplen $a_{i'}\le 2b_j$, así que el \textbf{mientras} avanza $i$ desde
$k_{j-1}+1\le k_j+1$ hasta llegar a $k_j+1$, y ahí para: o bien $k_j=p$ y $i=p+1$, o bien
$a_{k_j+1}>2b_j$. Entonces $p-i+1=p-k_j$ es justo la cantidad de $a_i>2b_j$, y se suma a
$cnt$. Al terminar el \textbf{para},
$cnt=\sum_{j=1}^q(p-k_j)=\sum_{j=1}^q\sum_{i=1}^p\ind{a_i>2b_j}=r(A,B)$.
\end{proof}

\begin{proposicion}
Para toda secuencia $L$ de largo $n\ge1$, la invocación \textsc{OrdenarYContar}$(L)$
devuelve $(R(L),L')$, donde $L'$ es la secuencia $L$ ordenada en forma creciente.
\end{proposicion}
\begin{proof}
Por inducción completa en $n$.

\emph{Base, $n=1$.} No hay pares $i<j$, así que $R(L)=0$, y una secuencia de un elemento
está ordenada. El algoritmo devuelve $(0,L)$.

\emph{Paso, $n\ge 2$.} Como $1\le m<n$ y $1\le n-m<n$, por HI
$(r_A,A')=(R(A),\text{$A$ ordenada})$ y $(r_B,B')=(R(B),\text{$B$ ordenada})$. Por el
Lema~\ref{l:cc}, $r=r(A',B')$, que por el Lema~\ref{l:inv} es igual a $r(A,B)$. Por la
correctitud del intercalado de MergeSort (sección 5.1 del libro), $L'=\textsc{Merge}(A',B')$
es la secuencia ordenada con los elementos de $A$ y $B$, o sea $L$ ordenada. Finalmente,
por el Lema~\ref{l:desc}, $r_A+r+r_B=R(A)+r(A,B)+R(B)=R(L)$.
\end{proof}

Notá que en ningún momento usamos que los números sean positivos: solo que $x\mapsto 2x$
es creciente, que es lo que garantiza $k_{j-1}\le k_j$.

\subsection*{(c) Tiempo de ejecución}

\begin{itemize}
\item \textsc{ContarCruzadas}$(A,B)$ con $|A|=p$, $|B|=q$: el \textbf{para} hace $q$
iteraciones con trabajo constante aparte del \textbf{mientras}. El cursor $i$ nunca
retrocede y solo avanza mientras $i\le p$, así que el cuerpo del \textbf{mientras} se
ejecuta a lo sumo $p$ veces \emph{en total}, sumando todas las iteraciones del
\textbf{para}. Las evaluaciones de la condición que dan falso son a lo sumo una por cada
$j$. El tiempo es $\Oh(p+q)$.
\item \textsc{Merge}$(A',B')$: cada paso agrega un elemento a la salida, así que hace
$p+q$ pasos de costo constante: $\Oh(p+q)$.
\item Dividir $L$ en $A$ y $B$ es $\Oh(n)$ (o $\Oh(1)$ si se trabaja con índices sobre un
arreglo).
\end{itemize}
Entonces existe una constante $c$ tal que el tiempo $T(n)$ de
$\textsc{OrdenarYContar}$ sobre una secuencia de largo $n$ cumple
\[
T(1)\le c,\qquad T(n)\le T(\lfloor n/2\rfloor)+T(\lceil n/2\rceil)+cn\quad(n\ge 2):
\]
$T(1)$ es el caso base (una comparación y un retorno); los dos términos recursivos
corresponden a las dos llamadas sobre las mitades; el término $cn$ es el costo de
dividir, contar las cruzadas y el intercalado.

Para $n$ potencia de $2$, y tomando $c$ suficientemente grande para que además
$T(2)\le 2T(1)+2c\le c'$ con $c'=4c$, queda $T(n)\le 2T(n/2)+c'n$ para $n>2$ y
$T(2)\le c'$, que es la recurrencia (5.1). Por (5.2) (Teorema~\ref{t:52}), $T(n)$ es
$\Oh(n\log n)$.

Si $n$ no es potencia de $2$, repetimos el argumento por inducción directa del
Ejercicio~1. Por inducción completa, $T(n)\le c\,n\lceil\log_2 n\rceil+c\,n$ para todo
$n\ge1$: para $n=1$ es $T(1)\le c$; para $n\ge2$, sea $\ell=\lceil\log_2 n\rceil$; como
$\lceil\log_2\lceil n/2\rceil\rceil=\ell-1$ y $1\le\lfloor n/2\rfloor\le\lceil n/2\rceil$,
ambas mitades cumplen $\lceil\log_2(\cdot)\rceil\le \ell-1$, y
\[
T(n)\le c\lfloor n/2\rfloor(\ell-1)+c\lfloor n/2\rfloor+c\lceil n/2\rceil(\ell-1)
+c\lceil n/2\rceil+cn=cn(\ell-1)+cn+cn=cn\ell+cn.
\]
En ambos casos $T(n)$ es $\Oh(n\log n)$.

%======================================================================
\section{Ejercicio 3}

\begin{enunciado}
El algoritmo de ordenación llamado Quicksort utiliza un método con el cual se
redistribuyen los elementos de un segmento de un arreglo en torno a un separador, también
llamado pivote, dejando el segmento parcialmente ordenado. Tras esa redistribución el
pivote queda en una posición, $s$, los elementos menores que el pivote quedan en las
posiciones anteriores a $s$, y los elementos mayores que el pivote quedan en las
posiciones posteriores a $s$. [\dots] Además el pivote queda exactamente en la posición que
quedaría si el segmento se ordenara totalmente.

El algoritmo \textsc{Particion} (figura 3.1) es una posible versión del método, en el cual
se elige el primer elemento del segmento como pivote. El tiempo de ejecución es
$\Theta(n)$, donde $n$ es la cantidad de elementos del segmento. Es un algoritmo
\emph{in place}: el espacio adicional usado es $\Oh(1)$.

\begin{algorithmic}[1]
\Require Un segmento de arreglo $A[\mathit{inf}..\mathit{sup}]$, $\mathit{inf}\le\mathit{sup}$
\Ensure El pivote es $A[\mathit{inf}]$; los elementos se redistribuyen: el pivote queda en
la posición $s$; cada elemento en $A[\mathit{inf}..s]$ es menor o igual al pivote; cada
elemento en $A[s+1..\mathit{sup}]$ es mayor que el pivote; devuelve $s$.
\State $pivote\gets A[\mathit{inf}]$;\ $s\gets\mathit{inf}$;\ $i\gets\mathit{sup}$
\While{$i>s$}
  \If{$A[i]\le pivote$}
    \State $s\gets s+1$; intercambiar $A[s]$ y $A[i]$
  \Else
    \State $i\gets i-1$
  \EndIf
\EndWhile
\State intercambiar $A[\mathit{inf}]$ y $A[s]$; \Return $s$
\end{algorithmic}

En cada iteración del ciclo la distancia entre $s$ e $i$ disminuye en uno, por lo que el
ciclo termina en $\Theta(n)$ iteraciones. En lo que sigue asumimos que todos los
elementos del arreglo son diferentes.
\begin{enumerate}[label=(\alph*)]
\item Considere la siguiente versión de Quicksort:
\textsc{QuickSort}$(A,\mathit{inf},\mathit{sup})$: si $\mathit{inf}<\mathit{sup}$ entonces
$s\gets\textsc{Particion}(A,\mathit{inf},\mathit{sup})$;
$\textsc{QuickSort}(A,\mathit{inf},s-1)$; $\textsc{QuickSort}(A,s+1,\mathit{sup})$.
Demuestre, resolviendo la relación de recurrencia, que el peor caso de tiempo de
ejecución es $\Theta(n^2)$. Describa una configuración de la entrada a la que le
corresponda el peor caso del tiempo de ejecución.
\item Asuma que se dispone del algoritmo \textsc{Mediana}, que dado un segmento de arreglo
devuelve la posición en que se encuentra la mediana de los elementos de ese rango (si la
cantidad de elementos es par, la que queda en la menor de las dos posiciones centrales). El
tiempo de ejecución de \textsc{Mediana} es $\Theta(n)$.
Usando \textsc{Particion} y \textsc{Mediana} diseñe un algoritmo que mediante la
estrategia dividir y conquistar ordene un arreglo en tiempo $\Oh(n\log n)$. Exprese el
tiempo de ejecución con una relación de recurrencia indicando la procedencia de cada
término. Explique por qué el tiempo de ejecución requerido es solución de esa relación de
recurrencia.
\end{enumerate}
\end{enunciado}

\examref{Examen diciembre 2012, Ej.~4 (QuickSort con pivote en la primera posición: peor
caso cuando la secuencia ya está ordenada, $\sum_{L=2}^{n}(L-1)=n(n-1)/2$
comparaciones); 2do.\ parcial 2013, parte (c) (QuickSort a partir de PARTITION, elementos de
dividir y conquistar). La idea de usar la mediana como separador para obtener
subproblemas de tamaño $\le n/2$ aparece en el Examen febrero 2025 (elemento faltante),
cuya solución oficial escribe la recurrencia indicando el origen de cada término y cita el
resultado del libro, igual que hacemos en la parte (b).}

\paragraph{Costo de \textsc{Particion}.} Sea $k=\mathit{sup}-\mathit{inf}+1$ el tamaño
del segmento. Al empezar, $i-s=k-1$; cada iteración disminuye $i-s$ en exactamente $1$
(o aumenta $s$, o disminuye $i$), y el ciclo termina cuando $i=s$. Por lo tanto hace
exactamente $k-1$ iteraciones, cada una con una comparación $A[i]\le pivote$ y costo
constante, más un trabajo constante fuera del ciclo. Entonces existen constantes
$0<a\le b$ tales que el tiempo de \textsc{Particion} sobre $k\ge1$ elementos, $P(k)$,
cumple $ak\le P(k)\le bk$.

\subsection*{(a) Peor caso de QuickSort: $\Theta(n^2)$}

Sea $W(n)$ el tiempo de peor caso de \textsc{QuickSort} sobre un segmento de $n$
elementos (distintos). Si $n\le 1$ (segmento vacío, $\mathit{inf}>\mathit{sup}$, o de un
elemento) solo se evalúa la condición: $W(0),W(1)\le d$ para una constante $d>0$. Si
$n\ge2$, \textsc{Particion} cuesta $P(n)$ y deja el pivote en la posición $s$; si el
pivote es el $(k+1)$-ésimo menor del segmento ($0\le k\le n-1$), la primera llamada
recursiva recibe los $k$ menores y la segunda los $n-1-k$ mayores. Entonces
\[
W(n)\le \max_{0\le k\le n-1}\bigl(W(k)+W(n-1-k)\bigr)+bn\qquad(n\ge 2).
\]
(El término $bn$ incluye la partición y el costo constante de la llamada.)

\paragraph{Cota superior: $W(n)\in\Oh(n^2)$.} Probamos por inducción completa que
$W(n)\le Cn^2+d$ para todo $n\ge0$, con $C=2b+d$. Para $n\le1$ es inmediato. Para $n\ge2$,
por HI (aplicable porque $k,\,n-1-k\le n-1$),
\[
W(n)\le \max_{0\le k\le n-1}\bigl(C k^2+C(n-1-k)^2\bigr)+2d+bn.
\]
La función $k\mapsto k^2+(n-1-k)^2$ es convexa, así que en el intervalo $[0,n-1]$ alcanza
su máximo en un extremo, donde vale $(n-1)^2$. Luego
\[
W(n)\le C(n-1)^2+2d+bn = Cn^2+d-\bigl[(2C-b)n-C-d\bigr].
\]
Como $2C-b=3b+2d>0$ y $n\ge2$, $(2C-b)n-C-d\ge 2(3b+2d)-(2b+d)-d=4b+2d>0$, y por lo tanto
$W(n)\le Cn^2+d$.

\paragraph{Cota inferior: una entrada que cuesta $\Omega(n^2)$.} Tomemos el arreglo
\textbf{ordenado en forma creciente}, $A[\mathit{inf}]<A[\mathit{inf}+1]<\dots<A[\mathit{sup}]$.
\begin{itemize}
\item En \textsc{Particion}, el pivote $A[\mathit{inf}]$ es el mínimo del segmento, así que
para todo $i>\mathit{inf}$ vale $A[i]>pivote$: en cada iteración se ejecuta $i\gets i-1$,
$s$ queda en $\mathit{inf}$, no se intercambia nada dentro del ciclo, y el intercambio
final es de $A[\mathit{inf}]$ consigo mismo. Devuelve $s=\mathit{inf}$ y el arreglo
\emph{no cambia}.
\item Entonces la primera llamada es sobre el segmento vacío
$A[\mathit{inf}..\mathit{inf}-1]$ y la segunda sobre $A[\mathit{inf}+1..\mathit{sup}]$,
de $n-1$ elementos, que \emph{sigue ordenado en forma creciente}.
\end{itemize}
Por lo tanto, si $T_{\mathrm{ord}}(n)$ es el tiempo con esta entrada,
\[
T_{\mathrm{ord}}(n)=T_{\mathrm{ord}}(n-1)+T_{\mathrm{ord}}(0)+P(n)\ \ge\ T_{\mathrm{ord}}(n-1)+an
\quad(n\ge2),\qquad T_{\mathrm{ord}}(1)>0.
\]
Resolviendo la recurrencia (desenrollándola):
\[
T_{\mathrm{ord}}(n)\ \ge\ T_{\mathrm{ord}}(1)+a\sum_{k=2}^{n}k\ \ge\
a\Bigl(\frac{n(n+1)}{2}-1\Bigr),
\]
que es $\Omega(n^2)$. Contando solo comparaciones con el pivote, el número exacto es
$\sum_{k=2}^{n}(k-1)=\frac{n(n-1)}{2}$, igual que en la solución oficial del examen de
diciembre 2012.

Como $T_{\mathrm{ord}}(n)\le W(n)\le Cn^2+d$, concluimos que $W(n)\in\Theta(n^2)$.

\paragraph{Configuración de peor caso.} En general el peor caso se da cuando en
\emph{todas} las llamadas el pivote (el primer elemento del segmento) es el mínimo o el
máximo del segmento, de modo que uno de los subproblemas queda vacío y el otro tiene
$n-1$ elementos; la cantidad de comparaciones con el pivote es entonces
$(n-1)+(n-2)+\dots+1=n(n-1)/2$. Ese es el máximo posible: si $K(n)$ es la máxima cantidad
de comparaciones, $K(0)=K(1)=0$ y $K(n)=\max_{0\le k\le n-1}(K(k)+K(n-1-k))+(n-1)$, y por
inducción completa $K(n)\le n(n-1)/2$, porque $k(k-1)/2+(n-1-k)(n-2-k)/2$ es convexa en
$k$ y en $[0,n-1]$ vale a lo sumo $(n-1)(n-2)/2$, y $(n-1)(n-2)/2+(n-1)=n(n-1)/2$. Un ejemplo concreto es el
\textbf{arreglo ya ordenado en forma creciente}, como vimos arriba.

(Notá que el arreglo ordenado en forma decreciente \emph{no} sirve como ejemplo con esta
versión de \textsc{Particion}: con $(5,4,3,2,1)$ la primera partición da $(2,1,4,3,5)$ y el
segmento izquierdo $(2,1,4,3)$ ya no tiene el pivote en un extremo.)

\subsection*{(b) QuickSort con la mediana como pivote: $\Oh(n\log n)$}

La idea es elegir como pivote la mediana del segmento. Como \textsc{Particion} usa como
pivote el primer elemento, antes de particionar se intercambia la mediana con
$A[\mathit{inf}]$.

\begin{algorithm}[H]
\caption{QuickSort con mediana}
\begin{algorithmic}[1]
\Require Un rango de un arreglo $A[\mathit{inf}..\mathit{sup}]$, con elementos distintos
\Ensure El rango queda ordenado en forma creciente
\Procedure{QuickSortMediana}{$A,\mathit{inf},\mathit{sup}$}
  \If{$\mathit{inf}<\mathit{sup}$} \Comment{si no, caso base: 0 o 1 elementos}
    \State $p\gets\Call{Mediana}{A,\mathit{inf},\mathit{sup}}$
    \State intercambiar $A[\mathit{inf}]$ y $A[p]$ \Comment{la mediana pasa a ser el pivote}
    \State $s\gets\Call{Particion}{A,\mathit{inf},\mathit{sup}}$ \Comment{dividir}
    \State \Call{QuickSortMediana}{$A,\mathit{inf},s-1$} \Comment{conquistar}
    \State \Call{QuickSortMediana}{$A,s+1,\mathit{sup}$}
  \EndIf \Comment{combinar: no hace falta, se trabaja in place}
\EndProcedure
\end{algorithmic}
\end{algorithm}
Para ordenar un arreglo $A[1..n]$ se invoca $\textsc{QuickSortMediana}(A,1,n)$.

\paragraph{Tamaño de los subproblemas.} Sea $n=\mathit{sup}-\mathit{inf}+1\ge2$. La mediana
es el elemento que ocupa la posición $\lceil n/2\rceil$ en el orden del segmento (si $n$ es
impar es la posición central $(n+1)/2$; si es par, la menor de las dos centrales, $n/2$).
Como los elementos son distintos, hay exactamente $\lceil n/2\rceil-1$ elementos menores y
$n-\lceil n/2\rceil=\lfloor n/2\rfloor$ mayores. Por la especificación de
\textsc{Particion}, el pivote queda en $s=\mathit{inf}+\lceil n/2\rceil-1$, y las
llamadas recursivas reciben
\[
n_1=\lceil n/2\rceil-1\le n/2\qquad\text{y}\qquad n_2=\lfloor n/2\rfloor\le n/2
\]
elementos, con $n_1+n_2=n-1$.

\paragraph{Correctitud.} Por inducción completa en $n$. Si $n\le1$ el rango ya está
ordenado y el algoritmo no hace nada. Si $n\ge2$, el intercambio deja en
$A[\mathit{inf}]$ a la mediana sin cambiar el multiconjunto del rango. \textsc{Particion}
deja el pivote en $s$, los elementos menores en $A[\mathit{inf}..s-1]$ y los mayores en
$A[s+1..\mathit{sup}]$ (los elementos son distintos, así que ``menor o igual'' equivale a
``menor'' salvo el propio pivote). Como $n_1,n_2<n$, por HI las dos llamadas ordenan sus
rangos; al reordenar cada rango se permutan elementos dentro de él, así que se mantiene que
todo elemento de $A[\mathit{inf}..s-1]$ es menor que $A[s]$, y este es menor que todo
elemento de $A[s+1..\mathit{sup}]$. Entonces $A[\mathit{inf}..\mathit{sup}]$ queda
ordenado. (Esta prueba es la misma que la de \textsc{QuickSort} de la parte (a); la
elección del pivote solo afecta el tiempo.)

\paragraph{Relación de recurrencia.} Sea $T(n)$ el tiempo de peor caso de
\textsc{QuickSortMediana} sobre $n$ elementos. Existe una constante $c>0$ tal que
\[
T(0),\,T(1)\le c,\qquad
T(n)\le T\bigl(\lceil n/2\rceil-1\bigr)+T\bigl(\lfloor n/2\rfloor\bigr)+cn\quad(n\ge2).
\]
Procedencia de cada término:
\begin{itemize}
\item $T(0),T(1)\le c$: caso base, solo se evalúa la condición $\mathit{inf}<\mathit{sup}$.
\item $T(\lceil n/2\rceil-1)$ y $T(\lfloor n/2\rfloor)$: las dos llamadas recursivas, sobre
los elementos menores y mayores que la mediana, cuyos tamaños calculamos arriba. Hay
\emph{dos} términos porque hay dos llamadas, y cada una tiene a lo sumo $n/2$ elementos
porque el pivote es la mediana.
\item $cn$: el trabajo no recursivo, que es \textsc{Mediana} ($\Theta(n)$), el intercambio
($\Oh(1)$), \textsc{Particion} ($\Theta(n)$, como vimos) y la comparación inicial
($\Oh(1)$). No hay costo de combinación.
\end{itemize}

\paragraph{Por qué la solución es $\Oh(n\log n)$.} Es más cómodo trabajar con
$U(m)=T(m-1)$, donde $m=n+1$ es la cantidad de elementos más uno. Con $m_1=n_1+1$ y
$m_2=n_2+1$ queda $m_1+m_2=n+1=m$ y $m_1,m_2\le n/2+1=(m+1)/2$, y la recurrencia es
\[
U(1),U(2)\le c,\qquad U(m)\le U(m_1)+U(m_2)+c\,m\quad(m\ge3).
\]

\emph{Caso $m$ potencia de 2} (o sea $n=2^k-1$). Entonces $n$ es impar, la mediana es el
elemento central y $n_1=n_2=(n-1)/2=2^{k-1}-1$, es decir $m_1=m_2=m/2$. Queda
$U(m)\le 2U(m/2)+cm$ para $m>2$ y $U(2)\le c$: es exactamente la recurrencia (5.1) del
libro, y por (5.2) (Teorema~\ref{t:52}) $U(m)$ es $\Oh(m\log m)$. Entonces $T(n)=U(n+1)$
es $\Oh(n\log n)$. Intuitivamente: en el árbol de recursión cada nivel tiene subproblemas
que suman a lo sumo $n$ elementos, y por lo tanto a lo sumo $cn$ (más el $+1$ por
subproblema) de trabajo, y como los tamaños se dividen por $2$ hay $\Oh(\log n)$ niveles.

\emph{Caso general.} Probamos por inducción completa que, para todo $m\ge1$,
\[
U(m)\le c\,m\log_2 m+3c\,m-2c .
\]
\textit{Base.} $U(1)\le c=0+3c-2c$; $U(2)\le c\le 2c+6c-2c$.
\textit{Paso} ($m\ge3$). Se cumple $1\le m_i\le (m+1)/2<m$, así que vale la HI para $m_1$
y $m_2$. Como $\log_2$ es creciente y $m_i\le(m+1)/2$,
$\log_2 m_i\le\log_2(m+1)-1$; y como $\ln(1+x)\le x$,
$\log_2(m+1)=\log_2 m+\log_2(1+1/m)\le\log_2 m+\frac{1}{m\ln 2}$. Entonces
\begin{align*}
U(m)&\le c\,(m_1\log_2 m_1+m_2\log_2 m_2)+3c(m_1+m_2)-4c+cm\\
&\le c\,m\bigl(\log_2(m+1)-1\bigr)+3cm-4c+cm\\
&\le c\,m\log_2 m+\frac{c}{\ln2}-cm+3cm-4c+cm
= c\,m\log_2 m+3cm-2c-\Bigl(2c-\frac{c}{\ln 2}\Bigr)\\
&\le c\,m\log_2 m+3cm-2c,
\end{align*}
porque $1/\ln 2\approx1{,}443<2$. Por lo tanto
$T(n)=U(n+1)\le c(n+1)\log_2(n+1)+3c(n+1)-2c$, que es $\Oh(n\log n)$.

\begin{observacion}
La cota $\Oh(n\log n)$ es la mejor posible: todo algoritmo de ordenación basado en
comparaciones necesita $\Omega(n\log n)$ comparaciones en el peor caso (cota del árbol de
decisión), así que el peor caso de \textsc{QuickSortMediana} es $\Theta(n\log n)$.
\end{observacion}

\end{document}
